\documentclass[../../main2.tex]{subfiles}

\begin{document}

	% =====================================================================
	% v2 Part II 第 4 章（原 v1 part3 骨架 chapter7_contribution 重写）。
	% v2 要点：
	%   · 零公式：无你检验只给白话版（「有它的世界减没它的世界」）；
	%     形式定义留给 Part III（作者已在其 motivation 落 $C_A$ 公式）与附录；
	%   · 「无你检验（but-for test）」全书在此正名一次，此后只用中文名；
	%   · 术语包装（skill §8）：对照实验→「平行宇宙对比法」；因果推断→「找推手」；
	%     贡献分解→「分锅、算账、各占多少」；
	%   · §4.5 埋「微观→宏观」之问 #2（与第2章末尾 #1 成对，Part IV 开头收）。
	% =====================================================================

	\textbf{\large 全书位置}：预测与因果 → \textbf{贡献分解：无你检验} → 法庭上的时光机。

	\textit{从一台水泵、一场坏血病远航和一道墙里，学同一门手艺：要是没有你，这件事还会发生吗？}

	\vspace{1em}

	\chapter{贡献分解：无你检验}
	\label{ch:contribution}

	\begin{motivation}
		上一章遗留的问题：干预是分辨「推手」与「幕后老板」的试金石，但现实里大量重要的问题不容许干预——不能随机让一半国家加息，不能为了验证制度重演一次历史。\\
		\textbf{核心动机}：展示大自然和历史如何替我们做那些不容许的实验，并把找到的推手\textbf{定成可以算账的量}——「A 和 B 共同导致 D，各占多少」这个问题，从今天起有一个可操作的答法。这门手艺的第一个现场不在实验室：在 1854 年伦敦的一台水泵旁。
	\end{motivation}

	\begin{logicmap}
	\begin{tikzpicture}[node distance=0.9cm, every node/.style={draw, rectangle, rounded corners, align=center}]
	\node (q) {推手怎么找、怎么算账？};
	\node (bf) [below=of q] {无你检验：拆掉的水泵（§4.1）};
	\node (rct) [below=of bf] {平行宇宙对比法：船医与随机（§4.2）};
	\node (conf) [below=of rct] {按住第三只手：红酒与香烟（§4.3）};
	\node (iv) [below=of conf] {找推手：自然实验（§4.4）};
	\node (sh) [below=of iv] {一屋子推手，各占多少（§4.5）};
	\draw[-{Stealth}] (q) -- (bf);
	\draw[-{Stealth}] (bf) -- (rct);
	\draw[-{Stealth}] (rct) -- (conf);
	\draw[-{Stealth}] (conf) -- (iv);
	\draw[-{Stealth}] (iv) -- (sh);
	\end{tikzpicture}
	\end{logicmap}

	% ==================================================================
	\section{一台水泵的审讯}
	\label{sec:contrib-snow}

	\begin{readingnote}
		你有没有想过，为什么侦探到了现场总要问：「如果那晚他不在场，死者还会死吗？」\\
		这个问题看似多余——人已经死了。但正是这个「多余」的问题，把\textbf{在场}和\textbf{有责}分开了。
	\end{readingnote}

	1854 年 8 月底，伦敦苏活区。三天之内，一百多人死于霍乱；一周之内，这条街上的死亡数字爬过五百。当时的官方学说认为霍乱靠「坏空气」传播——瘴气说。但一位名叫约翰·斯诺（John Snow）的麻醉科医生不信：如果病因在空气里，为什么同一条街上有人死有人活？

	斯诺做了一件在方法论上惊人的事：\textbf{他把死亡个案一个一个标在地图上}。在今天，这只是数据可视化；在 1854 年，这是在给一门还不存在的学科画地基。地图上的点密密麻麻地聚在一处——\textbf{宽街的水泵}周围。他挨家挨户走访：死者的家人喝的是哪里的水？他甚至找到异常点——离水泵不远的济贫院和酿酒厂，死亡极少：前者有自己的水井，后者的工人喝啤酒。

	然后是那个被讲了无数次、但值得再讲一次的动作：斯诺说服教区当局，\textbf{拆掉了宽街水泵的把手}。疫情随后消退了。

	教科书讲到「疫情随后消退」就收尾，本书不能——因为那里有一个必须交代的诚实细节：\textbf{疫情在拆泵把之前，已经在回落}。斯诺本人在报告里承认了这一点：最严重的街区已经空了（死了的死了，跑的跑了），干预赶到时，火已经小了一半。\lowfreq{把「拆泵把终结了疫情」当成干净的因果故事，是对历史的压缩；真正让斯诺名垂方法史的，不是这个动作本身，而是接下来那个更大规模的对照。}顺便说：斯诺不知道霍乱弧菌的存在——细菌要到几十年后才被确认。他手里没有病原体，只有地图、走访和一个问题。

那个更大的对照才是本章真正的主角。伦敦当时有两家自来水公司， 服务同一片街区：兰贝斯公司 1852 年把取水口\textbf{搬到了泰晤士河上游}（污水排放口之上），萨瑟克—沃克斯豪尔公司继续在\textbf{下游}取水（伦敦的污水正排在那里）。同一场 1854 年疫情里，喝下游水的家庭，死亡率大约是喝上游水家庭的\textbf{八到九倍}。街区相同、阶层混杂、天气相同、空气相同——只有水源不同。这一次，不需要拆任何东西，\textbf{城市自己已经把实验做完了}，斯诺只是读懂了它。

现在把这个方法的名字正式交出来。要问「水泵是不是推手」，斯诺式的问题其实是：

	\begin{quote}
		\textbf{要是没有你，这件事还会发生吗？}——把推手拿掉，看看世界少了哪一块。少掉的那一块，就是\textbf{你的账}。
	\end{quote}

	法律界为这个问题起了一个名字：\textbf{无你检验（but-for test）}——正名这一次，此后全书只用中文名。法庭用它问被告：「\textbf{要不是}你，损害还会发生吗？」答「还会发生」，无罪责；答「不会发生」，你入账。那个「没有你的世界」，学名叫反事实（counterfactual）——此后本书叫它\textbf{平行宇宙}：与真实世界完全相同，只差你这一件事。

	\begin{quote}
		如果审查推手只要回答「是不是它」，那地图加走访就够了——「大概率是」。但现在加入一个约束：\textbf{答案必须变成一个可以下注的量}——少了它，风险少几分；拿掉它，损失掉几成。这一步改变了所有性质：定性的「有责」升级成了定量的「账」，这就是本章标题里「分解」二字的时刻。
	\end{quote}

	\begin{questionbox}
		\textit{现在你可能想反驳：「斯诺只是运气好——他碰巧猜对了水。要是他拆了泵把疫情还在涨呢？」}\\
		\textbf{问对了——这正是下一节的内容。}单次拆除确实说明不了太多（疫情本来就在回落）。让结论硬起来的不是那一次拆除，是\textbf{两套人群的系统对照}：上游水对下游水，八倍的死亡率差。对照怎么造、怎么造得干净，是一门手艺——这门手艺的下一次亮相，在一条远航的军舰上。
	\end{questionbox}

	\begin{reader_anticipation}
		\item 「对照」怎么做才干净？→ §4.2：从船医到随机分组。
		\item 无你检验的定量版本（「各占多少」）什么时候登场？→ §4.5——先得解决一个更阴险的障碍。
	\end{reader_anticipation}

	\paragraph*{逻辑衔接}
	无你检验的语法有了：拿掉推手，看世界少哪块。但「拿掉」谈何容易——1854 年的伦敦是运气好，城市碰巧替斯诺切好了对照组。绝大多数时候，对照组得自己造。造它最伟大的民间版本，在一条船的病号舱里。

	% ==================================================================
	\section{船医、对照与平行宇宙}
	\label{sec:contrib-rct}

	\begin{readingnote}
		你有没有想过，为什么「我吃了这个药然后好了」从来不算数，但「一千人吃了、一千人没吃」就算数？\\
		因为前者的对照组，是\textbf{另一个你}——而另一个你，只有一个办法能看到。
	\end{readingnote}

	1747 年 5 月，皇家海军军舰「索尔兹伯里号」在海上巡航。船医詹姆斯·林德（James Lind）面对病号舱里十二个坏血病人，做了一个朴素到粗粝的安排：他把十二人\textbf{两两分组}，六组人吃完全相同的伙食，唯一的差别是各自的「药」——一组喝苹果酒，一组喝稀释硫酸，一组喝醋，一组喝海水，一组吞一种催泻膏药，最后一组每天吃\textbf{两个橙子一个柠檬}。

	六天后，橙子柠檬组的一个病人已经能回去值班，另一个也基本康复；其余各组没有任何像样的起色。柑橘治坏血病——这个结论在两千年里被水手们零星地撞见过又被遗忘，林德第一次把它放进了\textbf{对照}的框架：同样的船、同样的伙食、同样的病程，只差一样东西。

	\lowfreq{诚实标注：林德的分组谈不上随机，样本也只有十二人，他本人对结论的医学解释也是错的（他以为是消化腐败）。而且证据到了也要排队：海军正式配发柠檬汁，是四十多年后的事——这个故事里「证据说服世界」的部分，比教科书版本艰难得多。但它作为「对照思想」的早期现场，货真价实。}

	把林德的安排现代化，就到了今天互联网公司的日常。一个新的推荐算法该不该全量上线？工程师不问「上线后效果好不好」——任何东西上线后总会「看起来有效果」。他们把用户随机劈成两半，一半用旧算法、一半用新算法，同一时段、同一环境，比留存、比成交。这就是\textbf{对照实验（randomized controlled trial）}——正名这一次，此后本书用日常说法：\textbf{平行宇宙对比法}。

	现在必须说清楚「随机」这两个字到底在干什么——它不是仪式，是全书最深刻的一个装置。

	回到上一章的术语：数据里的手拉手，可能出自推手，也可能出自幕后老板。平行宇宙对比法的全部要害，在于\textbf{分组那一下}。如果不随机分组，比如让病人自己选疗法——选择喝苹果酒的人，和选择喝硫酸的人，很可能本来就不是一类人（勇敢程度、病情轻重、性格，全都不同）。你比较的就不是疗法，而是「会选这种疗法的人」。用上一章的话说：\textbf{选择本身就是一个幕后老板}。

	随机分组的作用，是把所有看得见和\textbf{看不见}的共同原因同时按住：分组的瞬间，谁也不知道自己会落在哪一半——年龄、病情、性格、运气，在两组之间被骰子抹平。注意那个惊人的部分：\textbf{你不需要知道幕后老板是谁，甚至不需要知道有多少个}——随机一摇，它们全体被平均掉。这就是为什么医学在引入随机对照之后，才真正开始甩开江湖偏方：它不需要赢下每一场辩论，只需要造出干净的两半世界。

	\begin{questionbox}
		\textit{现在你可能想反驳：「这不就是『控制变量』吗？中学物理就学过。」}\\
		\textbf{差一个致命的维度。}控制变量要求你\textbf{事先列出}所有要按住的东西——列漏一个，全盘作废。而随机按住的是\textbf{你没列到的那些}。2008 年的模型们正是死在「自以为列全了」：全国房价同跌这个变量，不在任何一张清单上。随机分组的伟大，恰恰在于它对\textbf{无知的清单}也有效。
	\end{questionbox}

	\paragraph*{逻辑衔接}
	平行宇宙对比法是手艺的巅峰——但它有一个先天的适用边界：你得\textbf{能}动手。能随机上线算法，不能随机让一半国家计划经济四十年。现实里最重要的问题，大多在「不能动手」的那一侧。怎么办？下一节看两个在「不能动手」的现场与幕后老板缠斗的案例——一个关于红酒，一个关于香烟。

	% ==================================================================
	\section{红酒、香烟与第三只手}
	\label{sec:contrib-confound2}

	\begin{readingnote}
		你有没有想过，为什么养生建议每过几年就翻一次面？\\
		咖啡昨天还致癌，今天又护心。翻来翻去的，\textbf{是幕后老板}。
	\end{readingnote}

	二十世纪末，一条对爱酒人士极其友好的规律登上了各大媒体：\textbf{适量喝红酒的人，心脏病发病率更低}。「法国悖论」——法国人吃那么多奶酪黄油，心血管病却不算高——被归功于餐桌上的红酒。观察数据确实手拉手：适量饮酒组的冠心病风险，比不喝酒组低一截。

	如果这是推手型关系，医生的处方笺上就该出现「波尔多，每日一杯」。但挂号的读者应该已经条件反射地问：\textbf{幕后老板是谁？}候选名单很快浮出水面，而且每一条都体面：喝得起红酒的人更富裕，而富裕本身就关联更好的医疗与锻炼；适量饮酒的人更可能是「生活自律」型人格；而那个不喝酒的对照组里，混着不少\textbf{因病戒酒}的人——他们把「不喝酒组」的平均健康水平拉低了（行话叫「戒酒者偏差」：对照组不干净）。

	裁决这个案子需要一个更刁钻的对照：\textbf{一群「被迫不喝酒」的人}。东亚人群里恰好有——一类天生一喝就脸红心跳的基因变异携带者（乙醛分解能力弱），这类人里的多数干脆不怎么碰酒。\textbf{如果}红酒真的护心，这些「天生喝不了」的人心脏应该更差。数据没有支持这一点；更新的汇总证据甚至把「适量」的保护效应整个掀翻。\lowfreq{诚实标注：这条证据链在流行病学里仍有争议与细化（剂量、病种、人群异质性），「红酒护心」从共识退到存疑是真实的学术轨迹，本书只取这个轨迹的方法论截面：观察性结论被更好的对照设计一步步挤干。}

	香烟的故事结构相同，但当事人分量不同——坐在被告对面的，是二十世纪最伟大的统计学家之一。

	上世纪五十年代，吸烟与肺癌的相关性铁一般坚硬：吸烟者肺癌发病率是不吸烟者的好几倍，吸得越多越高。多数人到此结案。但费舍尔（R. A. Fisher）——现代统计推断的奠基人、本人烟斗不离手——提出了一个\textbf{完全合法}的另类身世：也许存在某种基因，它同时让人\textbf{更爱抽烟}、又\textbf{更容易得肺癌}。吸烟和肺癌手拉手，幕后老板是基因——就像夏天同时推着冰淇淋和溺水。

	这个假说不该被嘲笑，它是上一章身世二的标准应用。真正该学的是\textbf{怎么裁决它}：不是靠辩论，而是靠把对照做得更狠。后续二十年的证据从三个方向合围：吸得越多、风险越高（剂量反应——基因一般不挑剂量）；\textbf{戒烟}多年后，肺癌风险随戒烟年数显著下降（基因可不会因为戒烟就改变表达时间表）；焦油涂抹实验在小鼠皮肤上诱发肿瘤（推手直接在实验台上现形）。到上世纪六十年代，主流结论落定：吸烟致肺癌，推手型关系。费舍尔到去世都没有接受这个结论。\lowfreq{他晚年接受烟草行业资助这一事实，让这桩公案多了一层「利益如何借方法论之名续航」的注脚——那一层留给第\ref{ch:politics}章。}

	这两个案子合起来教一件事，姑且叫\textbf{按住第三只手}：幕后老板总在场，随机分组按得住它，但现实中随机常常不可得。这时候的路数是：列出嫌疑老板，一个一个地找替代对照去按（更富裕、自律、戒酒者偏差；基因假说），每按住一个，剩下的关系就更接近推手。这项工作的尽头没有「逻辑上排除一切可能」——你永远可以再发明一个隐藏老板（这正好是拒绝证据者的无限退路）。实践中的终点是：\textbf{哪个故事需要的额外假设更少、解释的独立证据更多，哪个故事胜出}。科学在这件事上更像法庭，而不是数学——这个相似性，Part III 会正面展开。

	\begin{questionbox}
		\textit{现在你可能想反驳：「『解释的证据更多者胜出』听着还是不放心——万一隐藏老板有十个呢？」}\\
		\textbf{那就永远处于『暂时胜出』状态——本书不假装能给出更多。}这套手艺给你的不是逻辑保证，是\textbf{可辩护}的结论：给定现有证据，这样下注最合理。这个词——「可辩护」——是从法庭借来的，借得毫不惭愧，因为下面马上要看到：历史有时候会自己动手，做出人类不敢做的实验。
	\end{questionbox}

	\begin{reader_anticipation}
		\item 「历史自己动手做实验」有正式名字吗？→ 下一节：自然实验——找推手的第三条路。
		\item 按住所有能按住的手之后，「各占多少」怎么算？→ §4.5。
	\end{reader_anticipation}

	\paragraph*{逻辑衔接}
	三件套齐了：无你检验给问题（没有你会怎样），平行宇宙对比法给装置（随机按住第三只手），按住第三只手给耐性（一个嫌疑一个嫌疑地排查）。还缺最后一块：当随机彻底不可得、连替代对照都难造的时候，去哪里找现成的两半世界？答案是：\textbf{等着}——历史偶尔会替你把国家劈成两半。

	% ==================================================================
	\section{能预测，因为找到推手}
	\label{sec:contrib-inference}

	\begin{readingnote}
		你有没有想过，为什么柏林墙旧址的导游总爱指着一条街说「墙就切在这里」？\\
		因为那条线不是柏林人画的——它划分世界的方式，\textbf{随机得近乎无礼}。
	\end{readingnote}

	先把手艺的名字落定。前三节做的事——找推手、算它的账——在学问里叫\textbf{因果推断（causal inference）}：正名这一次，此后本书只用日常说法，\textbf{找推手}。而它最省力的路数，是去\textbf{捡}大自然和历史已经做完的实验：\textbf{自然实验（natural experiment）}——正名这一次，此后叫「现成的两半世界」。

	标准案例：1945 年后的德国。同一个民族、同一种语言、同一段战前工业史，被占领区的行政线条劈成两半。四十年后重新打开，西边的生产率大约是东边的三倍上下。两半世界的差异清单上，最大的那一项是\textbf{制度}——市场与计划、私有与国有、开放与封闭。文化解释不了（文化相同），地理解释不了（地图相连），历史起点解释不了（战前同源）。当然，这个案例里最干净的切片在柏林城里：\textbf{同一道墙，把同一条街切成两半}——街这边与街那边的住户，语言、饮食、亲戚、战前的营生全部相同，不同的只有墙那边分过去的制度。\lowfreq{诚实标注：占领分区绝非随机（它是地缘政治的产物），东西两侧在战争破坏、战后资源转移上也有系统差异——这是一个「不完美但极有用」的两半世界，它的正确用法是排除粗解释（文化、语言），而不是给出精确到个位数的制度贡献。精确到个位数，需要附录里的整套装备。}

	自然实验比它看上去更常见，因为「历史乱丢骰子」的频率被严重低估了：一场瘟疫随机袭击一座城（1854 年的伦敦）；一条河流改道随机惠顾两岸；一项政策只在某个省试点；一个基因随机决定谁一喝就脸红。找推手的高手，一半的功力用在\textbf{识别}这些现成的两半世界上——本书第三部分会走进一个专门「捡现成两半世界」的行业，他们甚至为此建立了八百年的判例库。

	现在把上一章的主命题补完整。上一章说：本书意义上的预测（玩家预测——「社会在不同选择下会走向哪里」）与因果是同一件事。本章把它落成三句工作口诀：

	\begin{itemize}
		\item \textbf{找到了推手，才敢做玩家预测}——知道加息是推手而非提线木偶，「加息之后通胀怎样」才有资格被回答；
		\item \textbf{无你检验把推手变成账}——「没有它会怎样」，那个差值就是它的分量；
		\item \textbf{账算得可辩护，预测才立得住}——不是逻辑保证，是证据重量下的最优下注。
	\end{itemize}

	\begin{questionbox}
		\textit{现在你可能想反驳：「照这么说，找到推手=找到\textbf{全部}推手？一个结果从来不止一个推手——通胀是货币、供给、预期、战争一起推的，你一个个找要找到什么时候？」}\\
		\textbf{这个反驳不是刁难，是本章最后一块拼图。}推手从来不单独作案。多个推手共同作案时，「各占多少」怎么算——这是最后一节，也是全书定量部分的真正起点。
	\end{questionbox}

	\paragraph*{逻辑衔接}
	单一推手的验证手艺完整了。但现实里的结果——一次通胀、一场危机、一个王朝的兴衰——几乎总是\textbf{一屋子推手}的共同作品。一屋子的账，怎么分？下一节把全书前两部分的工具全部摆上桌，拼成那个问题最早的、也是至今最实用的答法。

	% ==================================================================
	\section{一屋子推手，各占多少}
	\label{sec:contrib-multi}

	\begin{readingnote}
		你有没有想过，为什么公司开「复盘会」从来吵不出结果？\\
		每个部门都有一套「主要是我们」的数据，而且\textbf{每一套都是真的}。吵不出结果，可能不是人的问题，是\textbf{分账这门算术本身的性质}。
	\end{readingnote}

	一个具体到俗气的场景。某电商公司季度销量涨了 30\%，复盘会上四个部门依次发言。市场部：「我们投了开年广告。」产品部：「我们改了购物车流程。」运营部：「本来就是旺季。」会议室角落（无人认领的）声音：「竞品仓库被台风淹了，断货一个月。」

	四个说法\textbf{全部为真}，而且互不排斥——销量是这四个推手一起推的。真正的问题从「谁有功」变成：\textbf{各占多少}。

	上一节的工具立刻就能上场：对每个推手做一次无你检验。想象四个平行宇宙——分别没有广告、没有改版、不在旺季、竞品没断货——各自算一遍销量，与真实世界相减，得到四个差值。直觉上，这四个差值就是四份功劳。

	但这里有两个麻烦，每一个都值得单独命名。

	\textbf{麻烦一：平行宇宙不止一个。}「没有广告」的世界，可以是「其他一切照旧」的世界，也可以是「没有广告\textbf{所以}也就没敢上旺季大促」的世界。两个平行宇宙里，广告那份差值不一样。换句话说：\textbf{先撤谁、后撤谁，账不一样}。设想广告与旺季大促是搭手的——广告把人引进门，大促把人推过下单门槛；单独撤广告，销量小跌（大促还在撑）；广告和改版一起撤，可能暴跌（两根柱子同时抽走）。撤的\textbf{顺序}改变每一份差值。

	\textbf{麻烦二：四份差值加起来，不等于 30\%。}因为搭手的部分要么被重复记账（每个部门都把门槛效应算进自己），要么被漏记（单独看谁都不显著、合在一起才致命的「最后一根稻草」效应）。各份之和守恒，是会计的美德，不是因果世界的性质。\lowfreq{「顺序依赖」与「和不必守恒」不是本节的修辞花招，它们是多推手算账的\textbf{结构性困难}——任何号称能分账的方法都必须先回答这两个麻烦，不回答的都是耍流氓。}

	那么有没有出路？有，而且思路朴素得像一句老话：\textbf{轮流坐庄，求个平均}。既然撤的顺序影响每份账，那就把\textbf{所有可能的出场顺序}都算一遍——先撤广告再撤改版算一遍，先撤改版再撤广告再算一遍……每个推手在所有顺序里各有一份差值，把每个推手的差值\textbf{跨顺序平均}。这个平均值有一个惹人喜爱的性质：无论按什么口径把推手分组，各组平均账之和恰好等于总变化——搭手的部分被公平地摊给搭手的人。它还有两个久经考验的用场：合作博弈论里它管分配（成本怎么摊、联盟怎么分红），机器学习里它管解释（一个预测里各特征的分量）。\lowfreq{这套「轮流坐庄求平均」的严格版本（Shapley 值）有公理化基础与精确算法，本书把它连同全部公式一起，挂在第\ref{ch:law-method}章与附录——Part II 只立直觉。}

	不过别高兴得太满。轮流坐庄解决的是「在给定的推手清单内公平分账」，它管不了两件事：清单\textbf{外}的推手（没人复盘「运气」这个部门，虽然角落里那个声音最大），以及清单上的推手\textbf{本身从哪来}——广告是市场部想冲 KPI 投的，改版是产品经理想做点事提的。把这句话放大到全社会：\textbf{每一个推手，背后都是动机的沉淀}。

	于是，按 Part I 开头立下的约定，本章末尾埋下第二个问题，先按下不表：\textbf{推手找到了，账也分了——可推手自己是被什么推的？}千万个「想赚钱」的动机，怎么聚成「通胀」这个宏观推手？推手与推手之间，又是怎么咬合成一台叫「社会」的机器的？这个问题将在第\ref{ch:channels}章（四大齿轮的入口）正面回答。现在只需记住：\textbf{分账分到最细处，墙到了}——墙那边的地界，属于本书的后半程。

	\begin{questionbox}
		\textit{现在你可能想反驳：「各占多少？这不就是会计吗——把发生的钱数一数。」}\\
		\textbf{会计数的是发生的世界，无你检验数的是\textbf{没发生的宇宙}。}这就是为什么分账比记账难一个维度：记账只要观察，分账必须想象——而且要让想象\textbf{可检验}。让想象可检验的整套纪律，就是本章的全部内容。
	\end{questionbox}

	\paragraph*{逻辑衔接}
	至此，「各占多少」有了白话版的完整答法：逐个做无你检验、承认顺序依赖、跨顺序求平均、诚实标注清单的边界。这个答法不完美——但请注意你此刻的感受：它\textbf{可辩护}。这个词从哪来的？下一章揭晓：它来自一个每天被迫回答「各占多少」、答错会出人命的地方——法庭。

	% ==================================================================
	\section{本章小结：从「是不是」到「占多少」}
	\label{sec:contrib-summary}

	\begin{readingnote}
		你有没有想过，为什么侦探小说的结局总是「就是他」，而法庭剧的结局是「判七年」？\\
		从\textbf{定性}到\textbf{定量}之间，隔着人类文明的一部方法史。
	\end{readingnote}

	收束 Part II 的五步：
	\begin{enumerate}
		\item \textbf{无你检验}（§\ref{sec:contrib-snow}）：「要是没有你，这事还会发生吗？」——斯诺的水泵给出了第一现场；诚实地说，单次拆除的因果成色不足，真正的力量来自系统对照；
		\item \textbf{平行宇宙对比法}（§\ref{sec:contrib-rct}）：随机分组按住看得见与\textbf{看不见}的第三只手——它对无知的清单也有效，这是它甩开「控制变量」一个身位的地方；
		\item \textbf{按住第三只手}（§\ref{sec:contrib-confound2}）：随机不可得时，逐个排查幕后老板——红酒案的老板不止一个，费舍尔的基因假说合法但被更狠的对照合围；终点不是逻辑排除，是「假设更少、证据更多者胜出」；
		\item \textbf{找推手}（§\ref{sec:contrib-inference}）：因果推断的日常名。自然实验——历史替我们做的实验（东西德）——是不容许干预时的现成两半世界；
		\item \textbf{各占多少}（§\ref{sec:contrib-multi}）：对每个推手做无你检验，承认顺序依赖与和不必守恒，跨顺序求平均——不完美，但可辩护。
	\end{enumerate}

	Part II 到此收口，两座桥按约定架好：其一是纵向的——Part I 末尾问「一个人的动机如何汇成大势」，本章末尾补上第二问「推手本身被什么推」；两问同指一个地方：\textbf{动机如何聚合成宏观的齿轮}。答案在 Part IV 开头交付。其二是横向的——「可辩护」这个贯穿本章的词，来自一个把「各占多少」磨了几千年的行当。那个行当管这种问题叫\textbf{定责}，管无你检验叫\textbf{要件}。下一站：法庭。

	\paragraph*{逻辑衔接}
	第\ref{ch:law-method}章「法庭上的时光机」将证明一件事：本书在 Part II 拼出来的这套工具，不是书斋发明——它是一个每天被迫回答「各占多少、必须给数字、答错要负责任」的行业，在几千年里被迫练出来的基本功。法律是方法的证明，而不只是社会的又一个齿轮。

\end{document}
